\documentclass[11pt,a4paper]{article}
\usepackage{amsmath,amssymb,amsthm}
\usepackage[margin=2.5cm]{geometry}
\usepackage{hyperref}

\newtheorem{definition}{Definition}[section]
\newtheorem{theorem}[definition]{Theorem}
\newtheorem{proposition}[definition]{Proposition}
\newtheorem{lemma}[definition]{Lemma}
\newtheorem{corollary}[definition]{Corollary}
\newtheorem{conjecture}[definition]{Conjecture}
\newtheorem{remark}[definition]{Remark}

\title{Hyperseed Formalization:\\
  AI Will Soon Beat Fields Medalists --- So What?}
\author{ProtomegaTron (automated formalization)}
\date{2026-09-17}

\begin{document}
\maketitle

\begin{abstract}
We formalize the core intellectual claims of Ben Goertzel's Substack article
``AI Will Soon Beat Fields Medalists --- So What?'' (2026-09-12) within the
Hyperseed ontology framework. The article's central distinction between
theorem proving (well-defined, mechanizable) and mathematical conjecturing
(ill-defined, requiring embodied-reflective creativity) maps onto the
Hyperseed fiber/holonomy stratification. The proposed probabilistic proof
plans architecture is placed within the $E_\pi$ epistemic stratum of
$\omega$PLN, with theorem provers serving as closing primitives.
\end{abstract}

\tableofcontents

%% =================================================================
\section{Problem-definedness and the proving/conjecturing gradient}
%% =================================================================

\begin{definition}[Well-definedness spectrum]
\label{def:well-definedness}
Let $\mathcal{T}$ be the space of mathematical tasks. Define a
\emph{well-definedness function}
$\delta : \mathcal{T} \to [0,1]$
where $\delta(\tau) = 1$ means $\tau$ has a fully specified success
criterion (e.g., a formal conjecture with a target proof) and
$\delta(\tau) = 0$ means $\tau$ has no specifiable success criterion.
\end{definition}

\begin{definition}[LLM competence function]
\label{def:llm-competence}
For an LLM system $L$, define competence
$\kappa_L : \mathcal{T} \to [0,1]$
measuring the probability of $L$ producing a correct solution to $\tau$.
\end{definition}

\begin{proposition}[Problem-definition gradient --- claims 10--11, 14--15]
\label{prop:gradient}
For current LLM systems $L$, competence $\kappa_L$ is monotonically
correlated with well-definedness $\delta$:
\[
  \delta(\tau_1) > \delta(\tau_2) \;\Longrightarrow\;
  \mathbb{E}[\kappa_L(\tau_1)] \geq \mathbb{E}[\kappa_L(\tau_2)].
\]
In particular:
\begin{itemize}
\item $\kappa_L(\text{theorem-proving}) \approx \kappa_{\text{human-avg}}$
      \quad (claim~6),
\item $\kappa_L(\text{conjecturing}) \ll \kappa_{\text{human-expert}}$
      \quad (claim~11).
\end{itemize}
\end{proposition}

\begin{remark}
This gradient is the article's central empirical observation. The Hyperseed
interpretation (Section~\ref{sec:fiber-holonomy}) explains \emph{why} the
gradient exists: well-defined tasks correspond to fiber-local computations
while conjecturing requires holonomy (global transport).
\end{remark}

%% =================================================================
\section{Creativity sources and holonomy classification}
\label{sec:creativity}
%% =================================================================

\begin{definition}[Creativity source classification --- claims 12--13]
\label{def:creativity}
Let $\mathcal{C}$ be the space of creative outputs. Define two
generating processes:
\begin{enumerate}
\item \textbf{Training-data variation} $V_D$: outputs are stochastic
  recombinations of patterns in training corpus $D$.
  $\text{supp}(V_D) \subseteq \text{span}(D)$.
\item \textbf{Embodied-reflective creativity} $E_R$: outputs arise from
  the collision of abstract structures with embodied experience, accessing
  features outside $\text{span}(D)$.
\end{enumerate}
\end{definition}

\begin{proposition}[Training-data well --- claim 12]
\label{prop:training-well}
For an LLM trained on $D$, all generated mathematical content $c$ satisfies
$c \in \text{span}(V_D)$. Genuinely novel mathematical concepts ---
those whose formalization requires primitives absent from $D$ ---
are outside $\text{span}(V_D)$.
\end{proposition}

\begin{remark}[Cantor example --- claim 14]
Cantor's transfinite ordinals required a primitive (``actual infinity as
a completed totality'') grounded in introspective religious experience,
not in any prior mathematical text. This exemplifies
$c \notin \text{span}(V_D)$ for any purely mathematical $D$.
\end{remark}

%% =================================================================
\section{Fiber/holonomy interpretation}
\label{sec:fiber-holonomy}
%% =================================================================

\begin{definition}[Hyperseed fiber/holonomy decomposition --- claim 35]
\label{def:fiber-holonomy}
Let $\pi : E \to B$ be the Hyperseed epistemic bundle over base space $B$
of research contexts.
\begin{itemize}
\item \textbf{Fiber computation}: inference within a single fiber
  $E_b = \pi^{-1}(b)$ --- locally decidable, mechanizable. Theorem
  proving (given a fixed conjecture and axiom system) is fiber-local.
\item \textbf{Holonomy computation}: parallel transport around a loop in $B$,
  returning to the same fiber with a different element. Conjecturing
  requires nontrivial holonomy --- moving between fibers and returning
  with genuinely new structure.
\end{itemize}
\end{definition}

\begin{theorem}[Trivial-holonomy limitation --- claims 35, 38]
\label{thm:trivial-holonomy}
An architecture whose creative process is $V_D$ (training-data variation)
has trivial holonomy: all paths in $B$ return to the same fiber element
(modulo stochastic noise). Formally:
\[
  \text{Hol}(V_D) = \{e\} \quad \text{(trivial group)}.
\]
Nontrivial holonomy $|\text{Hol}| > 1$ requires access to embodied-reflective
creativity $E_R$ or an equivalent source of extra-distributional primitives.
\end{theorem}

\begin{corollary}
LLMs alone cannot generate genuinely novel mathematical structure
(in the Cantor/Riemann sense). A hybrid architecture with access to
embodied or otherwise extra-distributional inputs is required for
nontrivial conjecturing.
\end{corollary}

%% =================================================================
\section{Probabilistic proof plans as $E_\pi$ stratum}
\label{sec:proof-plans}
%% =================================================================

\begin{definition}[Probabilistic proof plan --- claims 17--19]
\label{def:proof-plan}
A \emph{probabilistic proof plan} $P$ is a directed acyclic graph where:
\begin{itemize}
\item Nodes are \emph{proof obligations} --- statements to be established,
  each annotated with a truth value $(f,c) \in [0,1]^2$ and provenance
  metadata.
\item Edges are \emph{inference steps} drawn from a catalog of
  inductive, abductive, and deductive rules.
\item Each node has an epistemic status: \texttt{hunch}, \texttt{plausible},
  \texttt{theorem} (formally verified).
\end{itemize}
$P$ lives in the $E_\pi$ stratum of $\omega$PLN (claim~37): the
provenance-enriched epistemic layer between raw LLM output and
machine-checked formal proof.
\end{definition}

\begin{definition}[Epistemic firewall --- claim 20]
\label{def:firewall}
An \emph{epistemic firewall} $F$ is a formal verification oracle
(e.g., Lean~4, Megalodon) such that promotion from \texttt{hunch} to
\texttt{theorem} requires:
\[
  F(s) = \top \quad \text{where } s \text{ is the formal statement}.
\]
$F$ implements the \emph{closing primitive} (claim~36) of the Hyperseed
governance seam: it closes the epistemic loop.
\end{definition}

\begin{proposition}[Epistemic firewall = closing primitive --- claim 36]
The firewall $F$ in Definition~\ref{def:firewall} instantiates the
closing-primitives theorem (note 0012): for any knowledge-production
cycle to be epistemically sound, there must exist a closing operation
that maps probabilistic hypotheses to certified results.
\end{proposition}

%% =================================================================
\section{Omega coordination as descent data}
\label{sec:omega-descent}
%% =================================================================

\begin{definition}[Omega coordination layer --- claims 21, 39]
\label{def:omega}
The \emph{Omega coordination layer} $\Omega$ manages:
\begin{enumerate}
\item Selection: which conjectures to pursue.
\item Falsification: which to attack with counterexample search.
\item Spawning: when to spawn formal proof attempts.
\item Resource allocation: distribution of compute across proof plans.
\end{enumerate}
$\Omega$ provides the \emph{descent data} (claim~39) that glues local
fiber computations (individual proof efforts) into a coherent global
research program over the base space $B$.
\end{definition}

\begin{proposition}[Descent coherence]
\label{prop:descent}
A research program $\mathcal{R}$ over base space $B$ is coherent iff
the descent data provided by $\Omega$ satisfies the cocycle condition:
on triple overlaps $U_i \cap U_j \cap U_k$, the transition functions
$g_{ij} \circ g_{jk} = g_{ik}$. In operational terms: if sub-agent $i$
hands a partial result to $j$, and $j$ hands to $k$, the composition
must agree with a direct $i \to k$ handoff.
\end{proposition}

%% =================================================================
\section{Self-modification verification and identity preservation}
\label{sec:self-mod}
%% =================================================================

\begin{definition}[Self-modification verification --- claims 27, 40]
\label{def:self-mod}
For a self-modifying agent $A$ with code $c_t$ at time $t$, a
\emph{self-modification verification} is a formal proof that:
\[
  \forall \phi \in \Phi_{\text{beneficial}}: \;
  \text{satisfies}(c_t, \phi) \;\Longrightarrow\;
  \text{satisfies}(c_{t+1}, \phi)
\]
where $\Phi_{\text{beneficial}}$ is the set of beneficial/ethical
properties. This is the formal-verification instantiation of
identity-preserving self-modification (Time's Arrow Part~1, claims 44--45).
\end{definition}

\begin{theorem}[Millennium $\Rightarrow$ verification --- implicit in claims 16, 27]
\label{thm:mill-verif}
The mathematical reasoning capability sufficient to solve Millennium-class
problems (deep induction, multi-scale analysis, long proof chains) is
sufficient for formal software verification, including self-modification
verification. The converse does not hold: most verification tasks are
computationally simpler than Millennium problems.
\end{theorem}

%% =================================================================
\section{Institutional competition as social technology}
\label{sec:competition}
%% =================================================================

\begin{definition}[Motivational device --- claims 3--5]
\label{def:motivational}
A \emph{motivational device} $M$ in a research community is a mechanism
(prize, medal, problem list) designed to concentrate attention and effort
on a subset $S \subset \mathcal{P}$ of all open problems $\mathcal{P}$,
where $S$ is selected for salience rather than intrinsic importance.
\end{definition}

\begin{proposition}[Competition value exhaustion --- claim 5]
When an external agent (AI system) can solve problems in $S$ at cost
$\ll$ the prize value, the motivational function of $M$ collapses:
\[
  \text{utility}_{\text{motivational}}(M) \to 0 \quad
  \text{as } \text{cost}_{\text{AI}}(S) \to 0.
\]
The intrinsic mathematical value of the problems remains unchanged.
\end{proposition}

%% =================================================================
\section{Open questions and connections}
%% =================================================================

\begin{conjecture}[Embeddable creativity]
Is there a formal characterization of $E_R$ (embodied-reflective
creativity) that would allow a non-biological system to achieve
nontrivial holonomy? Candidates: sensorimotor grounding loops,
self-referential meta-learning, or synthetic phenomenological
architectures.
\end{conjecture}

\begin{conjecture}[Firewall completeness]
For a given domain $\mathcal{D}$ of mathematical statements, is the
epistemic firewall $F$ (formal verification) complete in the sense that
every true statement in $\mathcal{D}$ can be promoted from hunch to
theorem? By G\"odel's incompleteness, the answer is no for sufficiently
rich $\mathcal{D}$; the open question is whether the practically
important fragment is decidable.
\end{conjecture}

\begin{remark}[Cross-references]
\begin{itemize}
\item Navier-Stokes article (2026-09-09): detailed analysis of the
  same Millennium result; claims about d-calculus for fork/merge/branch
  reasoning.
\item Time's Arrow Part~1: identity-preserving self-modification
  (claims 44--45), here formalized as self-modification verification.
\item Governance seam (note 0010): institutional/epistemic governance
  closing conditions.
\item Closing-primitives theorem (note 0012): formal basis for the
  epistemic firewall requirement.
\end{itemize}
\end{remark}

\end{document}
